\documentclass[header={Ancillary Note: Kernel And Range},notheorems]{study-note}

\begin{document}

\StudyTitleBlock{N4}{Ancillary Note}{Kernel, image, range, cokernel, and rank-type ideas}

\vspace{0.8em}

\begin{focusbox}
This note explains the standard algebraic objects attached to a linear map and why they matter for solving equations.
\end{focusbox}

\tableofcontents

\section{Core Definitions}

\begin{defbox}{Kernel}
For a linear map $T:X\to Y$,
\[
\ker T:=\{x\in X:Tx=0\}.
\]
\end{defbox}

\begin{defbox}{Image and range}
For a linear map $T:X\to Y$,
\[
\operatorname{Im}(T)=\operatorname{Ran}(T):=T(X).
\]
In this note, image and range mean the same subset of $Y$.
\end{defbox}

\begin{defbox}{Cokernel}
The \textbf{cokernel} of $T:X\to Y$ is the quotient space
\[
\operatorname{coker}(T):=Y/\operatorname{Ran}(T).
\]
It measures how far $T$ is from being surjective.
\end{defbox}

\section{Why They Matter}

\begin{focusbox}
For the equation
\[
Tx=y,
\]
the kernel controls uniqueness, and the range controls existence.
\end{focusbox}

\begin{thmbox}{Injectivity and surjectivity}
For a linear map $T:X\to Y$:
\begin{enumerate}
\item $T$ is injective if and only if $\ker T=\{0\}$;
\item $T$ is surjective if and only if $\operatorname{Ran}(T)=Y$;
\item $T$ is bijective if and only if both hold.
\end{enumerate}
\end{thmbox}

\begin{exbox}{Interpretation of the cokernel}
If $\operatorname{Ran}(T)\neq Y$, then some right-hand sides are not reachable. The quotient space
\[
Y/\operatorname{Ran}(T)
\]
records exactly the missing directions.
\end{exbox}

\section{Finite-Dimensional Picture}

\begin{thmbox}{Rank-nullity theorem}
If $X$ is finite-dimensional and $T:X\to Y$ is linear, then
\[
\dim X=\dim(\ker T)+\dim(\operatorname{Ran}(T)).
\]
\end{thmbox}

\begin{focusbox}
In finite dimensions, loss of dimension splits into two parts:
\begin{itemize}
\item collapse inside the domain, measured by the kernel;
\item missing directions in the codomain, measured by failure of surjectivity.
\end{itemize}
\end{focusbox}

\section{Functional-Analysis Context}

\begin{warningbox}{Closed range matters}
In Banach and Hilbert spaces, the range of a bounded operator need not be closed. This is one reason infinite-dimensional operator theory is subtler than finite-dimensional linear algebra.
\end{warningbox}

\begin{exbox}{Fredholm viewpoint}
In Fredholm theory one studies operators for which the kernel and cokernel are finite-dimensional. Then the defects of injectivity and surjectivity are both finite-dimensional and can be compared using the index.
\end{exbox}

\section{What To Remember Quickly}

\begin{focusbox}
Quick summary:
\begin{enumerate}
\item kernel = vectors sent to zero;
\item range = vectors that are actually hit;
\item kernel controls uniqueness;
\item range controls existence;
\item cokernel measures failure of surjectivity.
\end{enumerate}
\end{focusbox}

\end{document}
